% Full-population reduction and primitive-numerator obstruction.
% Independently audited analytic proofs; original mu remains unclassified.
\section{A population-only critical hierarchy and its primitive arithmetic}
\label{pop:section}

We use the original inert newborn with one real core, one real leaf,
and one inert parent stub. Let $n$ denote the real population and $d$
the retained root's full degree, including its stub. Every reachable
state satisfies $2\le d\le n$. Its only null event is the root choosing
the stub, of rate $1/d$. Thus the total proper-birth rate is
\[
 \lambda(T)=n-1/d.
\]
If a specified proper outcome has rate $b(T,T')$, its discounted
transition weight after all null events are summed is
\[
 \frac{b(T,T')}{s+\lambda(T)}.
\]
For an individual source-neighbour choice with source degree $d_u$,
this is $d/[d_u(ds+dn-1)]$. This is the same exact null-event
elimination already used by the finite profile-row calculator; here
it is applied to complete generation coefficients with a population
cutoff.

Let $L^*_{m,N}(s)$ count all depth-$m$ core-birth rewards whose prefix
has at most $N$ proper births, with no cutoff on null rings. It is a
finite rational function: each proper birth raises population by one.
For $s>0$,
\[
 L_{m,R,N}(s)\le L^*_{m,N}(s)\uparrow A_m^\infty(s)
 \quad(N\to\infty).
\]
The limit follows from monotone convergence, since every fixed birth
prefix has finite population and finitely many null events almost
surely. The root-preserving state recursion is explicit. If a core
$v$ has $\ell_v$ real leaves and full degree $d_v$, the rate of adding
a leaf at it is
\[
 \ell_v+\sum_{u\sim_Hv}\frac1{d_u},
\]
and its promotion rate is $\ell_v/d_v$. Divide each rate by
$s+n-1/d$, add the generation reward at a promotion, and iterate
through populations $2,\ldots,N+1$.

\subsection{Exact reduced pole denominators}

For $1\le m\le N$, define a set $\mathcal S_{m,N}$ of pairs $(n,d)$.
For $m\ge2$, put
\[
 \mathcal S_{m,N}
 =\{(n,d):2\le n\le N+1,\quad
                 2\le d\le\min(n,N-m+2)\}.
\]
For $m=1,N\ge2$, take all $2\le n\le N+1$, $2\le d\le n$,
except $(N+1,2)$. Finally $\mathcal S_{1,1}=\{(2,2)\}$. Set
\[
 B_{m,N}(s)=\prod_{(n,d)\in\mathcal S_{m,N}}(ds+dn-1),
 \qquad D_{m,N}=|\mathcal S_{m,N}|.
\]
Every displayed linear factor is primitive, since $\gcd(d,dn-1)=1$.

\begin{theorem}[No cancellation of any reachable pole]
\label{pop:poles}
The rational function $L^*_{m,N}$ has exactly the simple poles
\[
 s=-(n-1/d),\qquad(n,d)\in\mathcal S_{m,N}.
\]
At such a pole the residue has sign $(-1)^{n-2}$ and is nonzero.
Consequently $B_{m,N}$ is its primitive polynomial denominator,
up to an overall rational scalar, and it is also the primitive
polynomial denominator of $12(s+3)-L^*_{m,N}(s)$. Its degree is
\[
 D_{m,N}=
 \begin{cases}
 1,&m=N=1,\\
 N(N+1)/2-1,&m=1,\ N\ge2,\\
 [N(N+1)-m(m-1)]/2,&2\le m\le N.
 \end{cases}
\]
\end{theorem}
\begin{proof}
Expand the finite sum over proper-birth histories, keeping different
source choices separate. A reward prefix traverses successive
populations $2,3,\ldots,q$ and has a positive rational numerator
divided by $\prod_{j=2}^{q}(s+j-1/d_j)$. Different populations
have disjoint possible pole intervals $[-j,-j+1/2]$, and different
degrees at the same population give different poles. Each term
therefore has only simple poles.

At $s=-(n-1/d)$, all denominator factors from earlier populations
are negative, while all factors from later populations are positive.
There are exactly $n-2$ earlier factors. Thus every history using
this pole has residue of sign $(-1)^{n-2}$. The other histories are
analytic there. There can be no cancellation among its nonzero
residues. This argument evaluates the rational continuations of the
finite history weights, not divergent negative-discount expectations.

It remains to determine which poles occur. A depth-$m$ birth uses
at least $m$ promotions. Reaching root degree $d$ uses $d-2$ births
attached directly to the root, disjoint from those promotions.
Hence a retained reward history requires $m+d-2\le N$, and a
holding state before the reward requires $n\le N+1$.

For $m\ge2$ these conditions suffice. First grow the root's degree
to $d$. If $n\le d+m-1$, promote a chain for $n-d$ steps to reach
population $n$, then finish the chain to depth $m$. If
$n>d+m-1$, first build a depth-$(m-1)$ chain, then add the required
extra leaves at its last nonroot core, preserving a leaf there, and
finally promote a depth-$m$ child. The required birth count is at
most $N$ in both cases.

For $m=1$, a degree-$2$ root at population $n>2$ has only a core
neighbour and no real leaf. Its next depth-$1$ birth requires first
adding a root leaf and then promoting it, so it requires two more
births. This excludes precisely $(N+1,2)$ when $N\ge2$. Every
other pair is possible: for $d\ge3$, retain a root leaf while
growing a nonroot branch as necessary; for $d=2,n\le N$, grow
that branch and then add and promote a root leaf. The initial
$(2,2)$ state always occurs. Counting the stated pairs proves the
degree formulas. Subtracting $L^*_{m,N}$ from the polynomial
$12(s+3)$ reverses its nonzero residues and therefore leaves the
same reduced pole denominator.
\end{proof}

This is a denominator statement in $\mathbb Q(s)$. It does not
assert a lower bound on the reduced numerical denominator after
specializing $s$ to a rational number; further integer cancellation
at a specialization is a separate question.

\subsection{Integer forms after all polynomial cancellation}

Write $B=B_{m,N}$, and let $\kappa=\kappa_{m,N}$ be the least
positive integer such that
\[
 Q(s)=\kappa B(s)L^*_{m,N}(s)\in\mathbb Z[s].
\]
The integer exists, and
\[
 \kappa\mid K_N:=\prod_{n=2}^{N+1}\operatorname{lcm}(1,\ldots,n).
\]
Indeed, each source degree at population $n$ divides the corresponding
least common multiple, while the factor $d$ from null-event elimination
is a numerator. Multiplying a history by $K_NB$ leaves a polynomial
with nonnegative integer coefficients. Therefore $Q$ also has
nonnegative integer coefficients, $\deg Q\le D_{m,N}-m$, and
$Q(1)\le\kappa B(1)$.

Set
\[
 P(s)=12(s+3)\kappa B(s)-Q(s),\qquad
 g=\operatorname{content}(P),\qquad p_{m,N}=P/g.
\]
Here the content is positive, so $p_{m,N}$ is the primitive integer
numerator with positive leading coefficient. By Theorem~\ref{pop:poles},
no polynomial factor shared with $B$ remains. We have
\[
 \deg p_{m,N}=D_{m,N}+1,\qquad
 H(p_{m,N})\le49\kappa B(1)\le49K_NB(1).
\]
The height estimate follows as for the earlier certificates from the
nonnegative coefficients of $B,Q$ and $Q(1)\le\kappa B(1)$.

There is a useful bound on the possible integer content:
\begin{equation}\label{pop:content}
 \gcd(g,\kappa)=1,
 \qquad g\mid12\operatorname{lc}(B).
\end{equation}
For a prime dividing $\kappa$, its minimality ensures that some
coefficient of $Q$ is not divisible by that prime. Since
$P\equiv-Q$ modulo that prime, it does not divide $g$. On the
other hand the leading coefficient of $P$ is
$12\kappa\operatorname{lc}(B)$, proving the second assertion.

\subsection{A uniform obstruction even after primitive normalization}

\begin{theorem}[Primitive certificate values are never small]
\label{pop:primitive-growth}
For the original critical exponent $0<\mu<1$, every fully reduced
primitive numerator above satisfies
\[
 \boxed{p_{m,N}(\mu)>1\qquad(1\le m\le N).}
\]
More precisely, uniformly over $1\le m\le N$,
\[
 \log p_{m,N}(\mu)\ge N\log N-O(N),
\]
and the diagonal $m=N$ gives the matching upper order. Hence
\[
 \boxed{\min_{1\le m\le N}\log p_{m,N}(\mu)
             =N\log N+O(N).}
\]
These assertions concern this precise finite-population certificate
family; they do not exclude other signed arithmetic constructions.
\end{theorem}
\begin{proof}
Put $k=N+1-m$. Grow $k$ extra root leaves, then promote any root
leaf and make $m-1$ successive promotions along that new chain.
Its final depth-$m$ reward is at proper birth $N+1$, so it is omitted.
Each selected outcome class has rate at least $1/2$ (the first
promotion aggregates all root leaves), and at population $j$ its proper holding
rate is at most $j$. Its discounted mass is at least
\[
 \eta_{m,N}(s)\ge\frac{2^{-N-1}}{(s+2)_{N+1}}.
\]
An exact expression, if desired, is obtained with $d=k+2$:
\[
 \eta_{m,N}(s)=
 \left[\prod_{j=2}^{k+1}
       \frac{j(j-1)}{j(s+j)-1}\right]
 \frac{d-1}{d(s+d)-1}
 \prod_{r=1}^{m-1}\frac{d}{2[d(s+d+r)-1]}.
\]
At criticality, $A_m^\infty(\mu)\le12(\mu+3)$, so
$12(\mu+3)-L^*_{m,N}(\mu)\ge\eta_{m,N}(\mu)>0$.
Using~\eqref{pop:content},
\begin{equation}\label{pop:primitive-lower}
 p_{m,N}(\mu)
 \ge\frac{B(\mu)}{12\operatorname{lc}(B)}
          \frac{2^{-N-1}}{(\mu+2)_{N+1}}.
\end{equation}

Suppose first that $N\ge3$ and $m\le N-1$. For $m\ge2$,
the pole set includes $d=2$ at every population $2,\ldots,N+1$
and $d=3$ at every population $3,\ldots,N+1$. For $m=1$ the
same lower bound follows by replacing the missing final $d=2$
factor with the final $d=4$ factor; $N\ge3$ ensures its presence.
All other monic factors exceed one. Consequently
\[
 \frac{B(\mu)}{\operatorname{lc}(B)}
 \ge\frac32\prod_{n=3}^{N+1}(n-\tfrac12)^2.
\]
Since $\mu<1$, equation~\eqref{pop:primitive-lower} gives
\begin{equation}\label{pop:explicit-growth}
 p_{m,N}(\mu)\ge U_N:=
 \frac{(3/2)\prod_{n=3}^{N+1}(n-1/2)^2}
      {12\,2^{N+1}(3)_{N+1}}.
\end{equation}
Stirling's formula gives $\log U_N=N\log N-O(N)$.
Moreover
\[
 U_{10}=\frac{3454132356675}{2199023255552}>1,
 \qquad
 \frac{U_{N+1}}{U_N}=\frac{(N+3/2)^2}{2(N+4)}>1
 \quad(N\ge3).
\]

For $N\le9$, each history earns at most $N$ rewards and hence
$L^*_{m,N}(\mu)\le N$. The content bound directly gives
\[
 p_{m,N}(\mu)
 \ge\frac{B(\mu)}{12\operatorname{lc}(B)}(36-N)
 \ge\frac{27}{8}>1.
\]
Here the pole set is nonempty and every monic factor exceeds $3/2$.

Finally, when $m=N$, every retained depth-$N$ reward requires
exactly $N$ consecutive chain promotions and no other proper births.
The root degree stays two, so
\[
 L^*_{N,N}(s)=\frac1{\prod_{n=2}^{N+1}(2s+2n-1)},\qquad
 p_{N,N}(s)=12(s+3)\prod_{n=2}^{N+1}(2s+2n-1)-1.
\]
This numerator is primitive: its leading coefficient is $12\,2^N$,
whereas its constant coefficient is coprime to both two and three.
It exceeds one at every positive $s$, and
$\log p_{N,N}(\mu)=N\log N+O(N)$. This handles every diagonal
case and supplies the matching upper order for the minimum.
\end{proof}

\begin{corollary}[A stronger algebraic-root hierarchy]
\label{pop:roots}
Let $\zeta_{m,N}=0$ when $L^*_{m,N}(0)\le36$; otherwise let it
be the unique positive root of $L^*_{m,N}(s)=12(s+3)$. Then
\[
 \zeta_N=\max_{1\le m\le N}\zeta_{m,N}
 \quad\hbox{satisfies}\quad
 0\le\zeta_N\le\zeta_{N+1}<\mu,\qquad
 \zeta_N\longrightarrow\mu.
\]
It dominates the earlier diagonal-budget root $\xi_N$.
For $N\ge2$, it is the largest nonnegative real root of
\[
 s\prod_{m=1}^N p_{m,N}(s),
 \qquad
 \deg\left(s\prod_{m=1}^N p_{m,N}\right)
 =\frac{N(N+1)(2N+1)}6+N.
\]
For $N=1$ the degree is three. No effective convergence rate in $N$
for these critical roots is asserted.
\end{corollary}
\begin{proof}
Each lower sum is a finite positive Laplace transform and decreases
to zero with $s$. Its threshold residual is strictly increasing, and
is positive at $\mu$ by the omitted-history proof. Increasing $N$
retains every old reward, proving monotonicity. For every $s<\mu$,
the complete lower-certificate theorem supplies some $m,R,N$ with
$L_{m,R,N}(s)>12(s+3)$; domination by $L^*_{m,N}$ then proves
convergence to $\mu$. The same domination at equal diagonal budgets
gives $\zeta_N\ge\xi_N$. The primitive polynomials have the same
nonnegative roots as their threshold residuals. Summing their exact
degrees from Theorem~\ref{pop:poles} gives the product degree.
\end{proof}

The theorem shows that cancelling all spurious polynomial factors and
then dividing all integer content does not produce small primitive
critical values for this family. A different signed construction is
required for the stated small-value route to irrationality or transcendence.

\subsection{An explicit population-only error bound}
\label{pop:tail-section}

The proper-birth calculation also admits an outer error that depends
only on its population cutoff. Let $Y_m(s)$ be the total discounted
depth-$m$ reward in the potential genealogy in which every individual
has children at rate one forever. The actual core genealogy embeds
in this process by keeping every source clock but retaining only the
actual promotions. All assertions below use fixed $m$.

\begin{proposition}[Population-only error with an explicit rate]
\label{pop:tail}
For $s>0$, $m\ge1$, $N\ge1$, and $t>0$, put
\[
 M_m(s)=s^{-2m}\sum_{j=0}^{m}(s/2)^j.
\]
Then
\begin{align*}
0\le A_m^\infty(s)-L^*_{m,N}(s)
&\le s^{-m}e^{-st}\sum_{j=0}^{m-1}\frac{(st)^j}{j!}\\
&\quad+
\sqrt{M_m(s)(1-e^{-t})^N(1+Ne^{-t})}.
\end{align*}
In particular, for each fixed $s,m$,
\[
 A_m^\infty(s)-L^*_{m,N}(s)
 =O_{s,m}\bigl(N^{-s}(\log N)^{s+m-1}\bigr).
\]
This is an effective fixed-generation approximation. It supplies no
outer error for letting the generation index tend to infinity in the
definition of the pressure.
\end{proposition}
\begin{proof}
Let $\tau_N$ be the time of the $N$th actual proper birth. Every
omitted depth-$m$ reward is after $\tau_N$. Under the potential
genealogy coupling, its sum is bounded by
\[
 Y_m^{>t}(s)+Y_m(s)\mathbf1_{\{\tau_N<t\}},
\]
where $Y_m^{>t}$ counts potential births after time $t$.
The expected generation-$m$ birth intensity is
$u^{m-1}/(m-1)!$, so integration gives
\[
 \mathbb E Y_m^{>t}(s)
 =s^{-m}e^{-st}\sum_{j=0}^{m-1}\frac{(st)^j}{j!}.
\]
Poisson offspring decomposition gives
\[
 \mathbb E Y_m(s)=s^{-m},\qquad
 \mathbb E Y_m(s)^2
 =s^{-2m}+\frac{\mathbb E Y_{m-1}(s)^2}{2s}.
\]
Starting from $Y_0=1$, this solves to $M_m(s)$. These identities
follow first for finite time and generation truncations, with the
nonnegative limit justified by monotone convergence.

The actual proper-birth rate at population $n$ is at most $n$.
Therefore its real population is dominated by the Yule process from
two founders, and
\[
 \Pr(\tau_N<t)
 \le\Pr(Y_2(t)\ge N+2)
 =(1-e^{-t})^N(1+Ne^{-t}).
\]
Cauchy--Schwarz proves the displayed error bound; it requires no
independence between $\tau_N$ and the potential reward.

For the rate, choose $c=2s+2$ and
$e^{-t}=c\log N/N$ whenever this lies below one. The first term is
$O_{s,m}(N^{-s}(\log N)^{s+m-1})$, while
$(1-e^{-t})^N\le N^{-c}$ makes the second term
$O_{s,m}(N^{-s-1}\sqrt{\log N})$.
\end{proof}

At a rational discount $s=a/b$, the bound can be made entirely
rational. Choose $t=k\log2$, so $e^{-t}=2^{-k}$. Enclose
$2^{-ks}$ above by rational bisection of its $b$th-power equation,
use $st\le sk$ in the finite polynomial, and enclose the remaining
square root above by an integer square comparison. Thus no
floating-point inequality is needed for the outer error. The
implemented helper follows exactly this prescription.

\subsection{The remaining population cost}

Let
\[
 q_i(s)=\prod_{j=2}^{i+1}\frac{j}{j+s},\qquad
 V_N(s)=\sum_{i=1}^{N}q_i(s).
\]
The proper-birth waiting times can be coupled as
$E_i/\lambda_i\ge E_i/(i+1)$ with fresh independent unit
exponentials. There is at most one generation reward at each proper
birth, so
\[
 L^*_{m,N}(s)\le V_N(s).
\]
The same finite telescoping calculation as for the core-ring envelope
gives
\[
 V_N(s)=
 \begin{cases}
 \displaystyle\frac{(N+2)q_N(s)-2}{1-s},&s\ne1,\\[5pt]
 2H_{N+2}-3,&s=1.
 \end{cases}
\]
At $a=113/200$, exact rational arithmetic gives
\[
 42.761044675975962399<V_{487}(a)<42.761044675975962400,
\]
\[
 42.803125001232158295<V_{488}(a)<42.803125001232158296.
\]
Since $12(a+3)=42.78$, any threshold root from this population-only
family in the current enclosure requires $N\ge488$. This is only
necessary. Eliminating the core-ring budget does not make enumeration
at that population practical, and no new numerical bracket for $\mu$
is claimed.
